%\documentclass[wcp,gray]{jmlr} % test grayscale version
\documentclass[wcp]{jmlr}

% Automatically loaded: amsmath, amssymb, natbib, graphicx, url, algorithm2e
\usepackage{longtable}
\usepackage{booktabs}
\usepackage{multirow}
\usepackage{subcaption}
\usepackage{xcolor}
\usepackage{hyperref}

\definecolor{darkblue}{rgb}{0, 0, 0.5}
\hypersetup{colorlinks=true, citecolor=darkblue, linkcolor=darkblue, urlcolor=darkblue}

\pagenumbering{gobble}
\newcommand{\cs}[1]{\texttt{\char`\\#1}}
\makeatletter
\let\Ginclude@graphics\@org@Ginclude@graphics
\makeatother

\jmlrvolume{}
\jmlryear{2026}
\jmlrworkshop{ACML 2026}

\title[Prompt Elaboration: Supplementary Material]{When Does Prompt Elaboration
       Stop Helping? Supplementary Material}

% Anonymous for double-blind review
\author{\Name{Anonymous Author(s)} \Email{anonymous@example.com}\\
\addr Anonymous Institution}

\editors{TBD}

\begin{document}

\makeatletter
\let \@jmlrpages \@empty
\makeatother

\maketitle

\noindent
This supplementary material accompanies the main paper and provides the full
prompt specifications, the LLM judge rubric, the complete per-model saturation
statistics, and worked prompt walkthroughs. It is not required to follow the
main paper's claims, which are self-contained, but documents the exact prompts
and per-model results for reproducibility. Section, table, and figure numbers
here are independent of the main paper.

\appendix
\input{sec_supplementary}

\end{document}
